\chapter{Introdução}
\label{cap:introducao}

\section{Contexto e motivação}
\label{sec:contexto}

O presente trabalho, ainda em fase de desenvolvimento, tem como principal
objetivo explorar a utilização do conhecimento especializado, através da
troca de ideias com um perito e regras do domínio, no contexto da tomada de
decisão aplicada à condução autónoma. Uma vez que o foco do trabalho são os
sistemas periciais, há que ter em conta não só a base de conhecimento que
descreve factos, relações e regras do domínio, mas também o motor de
inferência que avalia esse conhecimento, determina as regras a utilizar e
permite desenvolver novas conclusões. Neste
sentido, o objetivo do sistema pericial vai para além de dar uma resposta ao
utilizador: é compreender e justificar as decisões tomadas.

O trabalho desenvolvido ao longo destas quatro semanas parte do princípio
referido anteriormente e foca-se em processar, em tempo real, os dados
provenientes de determinados sensores simulados. Procurou-se diferenciar o
trabalho desenvolvido através do uso de um domínio dinâmico, ou seja, em vez
de consultar o utilizador através de uma sequência estática de perguntas, o
sistema recebe continuamente informação proveniente dos sensores, transforma
essas observações em factos e mantém na \emph{Working Memory} o conhecimento
relevante para o estado atual do veículo e do ambiente envolvente. De
momento, encontra-se em implementação um motor de regras Drools, integrado
numa aplicação, a partir da qual é possível fazer a simulação de cenários e
processar os dados contínuos dos sensores. A utilização de sessões
\emph{stateful}, juntamente com o processamento temporal de eventos, permite
manter a informação relevante entre diferentes ciclos de decisão e, desta
forma, ter em conta a evolução do ambiente ao longo do tempo. A partir da
informação disponível, o motor de inferência avalia as regras cujas condições
se encontram satisfeitas e apresenta propostas de decisão que podem,
posteriormente, ser comparadas de acordo com os critérios de prioridade
definidos para o sistema.

Uma das principais dificuldades do projeto resulta da possibilidade de várias
regras serem válidas em simultâneo e conduzirem a propostas de ação
diferentes. Neste sentido, foram considerados mecanismos de prioridade e de
resolução de conflitos, que permitem definir uma ordem entre as diferentes
propostas, em que as situações de maior risco têm maior importância na
tomada de decisão. A explicabilidade das decisões é, igualmente, uma
componente central do sistema: para além de o sistema apresentar uma ação,
deve ser possível compreender por que razão essa ação foi escolhida, como foi
concluída e qual o motivo pelo qual as outras alternativas não foram
selecionadas.

No final do Desafio~I, espera-se que, com a abordagem descrita, seja possível
explorar alguns dos principais princípios associados aos sistemas periciais,
nos quais o conhecimento e o processo de inferência devem estar acessíveis
para análise e validação. No decorrer do desenvolvimento do trabalho, foram
consideradas abordagens focadas nas medidas temporais de segurança como, por
exemplo, \emph{Time To Collision} (TTC), bem como abordagens complementares,
entre as quais a Teoria dos Jogos, analisada como possibilidade de evolução
futura para determinadas situações de interação entre veículos. Importa ainda
referir que o objetivo do presente projeto não consiste em desenvolver ou
reproduzir um sistema completo de condução autónoma, mas sim em estudar o
processo de decisão e a forma como este pode ser estruturado, acompanhado,
analisado e justificado.

\section{Problemática e questões orientadoras}
\label{sec:problematica}

O problema abordado neste projeto diz respeito à forma como a informação
recolhida, continuamente, pelos sensores de um veículo autónomo pode ser
utilizada por um sistema pericial para apoiar a tomada de decisão durante a
condução. Ao longo do percurso, o veículo recebe diferentes tipos de
informação sobre o ambiente/cenário como, por exemplo, reconhecer a
velocidade atual, identificar um veículo à frente, detetar um peão ou
reconhecer determinadas condições meteorológicas. Contudo, esta informação
não é necessariamente recebida em simultâneo, uma vez que os diferentes
sensores têm alcances e frequências de atualização distintos.

Os dados recolhidos pelos sensores são transformados em eventos e factos que
representam aquilo que o sistema conhece sobre a situação atual. Alguns
desses factos podem manter-se durante várias leituras como, por exemplo, o
limite de velocidade, enquanto outros podem sofrer alterações entre
diferentes leituras, como a distância a um obstáculo ou o tempo estimado até
uma possível colisão. Com base nesta informação, o motor de inferência
verifica quais são as regras da base de conhecimento cujas condições são
satisfeitas em cada situação. Tendo em conta que mais do que uma regra pode
ser válida no mesmo ciclo de inferência, o sistema deve, neste sentido,
comparar as diferentes propostas e selecionar a ação que deverá prevalecer,
tendo em conta as prioridades definidas para a tomada de decisão.

Após a tomada de decisão, o sistema poderá determinar diferentes ações como,
por exemplo, manter a velocidade, desacelerar, travar ou realizar uma mudança
de via. Como essas ações podem alterar o estado do veículo e,
consequentemente, a situação observada, os sensores continuam a recolher nova
informação sobre o ambiente. Estes novos dados são novamente encaminhados
para o processo de inferência, para dar início a um novo ciclo de decisão. É
esta interação contínua que permite ao sistema acompanhar a evolução da
situação ao longo do tempo. Neste sentido, o funcionamento pretendido pode
ser representado, de uma forma simplificada, através do seguinte processo:

\begin{figure}[H]
  \centering
  \begin{tikzpicture}[
      node distance=7mm and 6mm,
      etapa/.style={draw=corEscura, fill=corFundo, thick, rounded corners=2pt,
        text width=2.9cm, minimum height=1.3cm, align=center, font=\small},
      seta/.style={-latex, thick, draw=corEscura}]
    \node[etapa] (sensores) {Leitura dos\\sensores};
    \node[etapa, right=of sensores] (factos) {Eventos e factos na memória de trabalho};
    \node[etapa, right=of factos] (motor) {Motor de\\inferência:\\avaliação das regras};
    \node[etapa, right=of motor] (propostas) {Propostas de\\decisão};
    \node[etapa, below=of propostas] (selecao) {Seleção da decisão segundo as prioridades};
    \node[etapa, left=of selecao] (acao) {Ação sobre o\\veículo};
    \node[etapa, left=of acao] (estado) {Novo estado do veículo e do ambiente};
    \draw[seta] (sensores) -- (factos);
    \draw[seta] (factos) -- (motor);
    \draw[seta] (motor) -- (propostas);
    \draw[seta] (propostas) -- (selecao);
    \draw[seta] (selecao) -- (acao);
    \draw[seta] (acao) -- (estado);
    \draw[seta] (estado.west) -| node[pos=0.3, below, font=\footnotesize,
      text=corCinza] {nova leitura} (sensores.south);
  \end{tikzpicture}
  \caption{Processo de inferência e tomada de decisão do sistema pericial}
  \label{fig:ciclo-decisao}
\end{figure}

A partir do problema apresentado, foi definida a seguinte questão para
orientar o desenvolvimento do projeto:

\begin{quote}
  \itshape De que forma pode um motor de inferência utilizar, continuamente,
  dados provenientes de diferentes sensores para propor decisões na tomada de
  decisão, na condução, tendo em conta a evolução temporal da informação?
\end{quote}

Tendo em conta esta questão central, foram definidas algumas questões
orientadoras mais específicas, nomeadamente:
\begin{itemize}
  \item De que forma podemos utilizar os dados provenientes de diferentes
    sensores para manter uma representação atualizada da situação do veículo
    e do ambiente?
  \item Até que ponto a análise de várias leituras ao longo do tempo pode
    ajudar a distinguir uma situação pontual de uma situação de perigo
    persistente?
  \item Como pode o sistema selecionar uma decisão quando várias regras são
    aplicáveis em simultâneo e apresentam propostas de ação diferentes?
  \item Como pode o sistema explicar o processo que o levou a uma determinada
    conclusão ou ação, através da identificação de factos e respetivas regras
    (\emph{How})?
  \item Como pode o sistema explicar por que razão uma determinada decisão foi
    tomada e por que motivo não foram escolhidas outras alternativas
    (\emph{Why} e \emph{Why Not})?
  \item De que forma o conhecimento adquirido através da pesquisa e contacto
    com o perito pode ser transformado em factos e regras e, posteriormente,
    validado e melhorado?
\end{itemize}

\section{Estrutura do relatório}
\label{sec:estrutura}

Este relatório descreve o trabalho realizado até ao momento e distingue duas
camadas que convém não confundir: o modelo de conhecimento, construído a
partir da pesquisa e das sessões com o perito, que define o comportamento
pretendido para o sistema, e o protótipo atual, desenvolvido em Drools, que
implementa já uma parte desse modelo. Ao longo do texto, o que já existe no
protótipo é descrito no presente, enquanto o que está apenas previsto é
assinalado como tal, e o estado de cada mecanismo é resumido no
\cref{cap:implementacao}.

O relatório começa por apresentar os objetivos do projeto e o processo de
aquisição de conhecimento, incluindo as fontes utilizadas e as sessões com o
perito. Descreve depois os sensores considerados e a forma como as suas
leituras dão origem a factos, a representação desse conhecimento e o modo
como o sistema raciocina ao longo do tempo, resolve conflitos entre propostas
de decisão e explica as escolhas que faz. Seguem-se a arquitetura e a
implementação do protótipo, os cenários e critérios utilizados para o
validar e, por fim, as limitações, o trabalho futuro e as conclusões.
